\documentclass[12pt, ngerman, a4paper,toc=left, parskip=half, numbers=noenddot, headings=standardclasses, nonumberlist]{scrartcl}
\usepackage[utf8]{inputenc}
\usepackage[T1]{fontenc}
\usepackage{babel}
%\usepackage{lmodern}
\usepackage[top=2.5cm, right=2.5cm, bottom=2cm, left=2.5cm, footskip=0.5cm]{geometry}
\usepackage{calc}
\usepackage{graphicx}
\usepackage{amsmath}
\usepackage{tabularx}
\usepackage{booktabs}
\usepackage[official]{eurosym}
\usepackage{lastpage}
\usepackage{etoolbox}
\usepackage{amssymb}
\usepackage{enumitem}
\usepackage{natbib}
\bibliographystyle{agsm}
\usepackage[hidelinks]{hyperref}
\usepackage{dramatist}
\usepackage{listings}

\usepackage[table]{xcolor}
\usepackage{multirow}
\usepackage{array}

\usepackage{adjustbox,lipsum}
\usepackage{makecell}

\lstset{ 
    breakatwhitespace=false,         % sets if automatic breaks should only happen at whitespace
    breaklines=true,                 % sets automatic line breaking
    frame=single,                    % adds a frame around the code
    keepspaces=true,                 % keeps spaces in text, useful for keeping indentation of code (possibly needs columns=flexible)
    numbers=left,                    % where to put the line-numbers; possible values are (none, left, right)
    numbersep=5pt,                   % how far the line-numbers are from the code
    numberstyle=\color{gray},        % the style that is used for the line-numbers
    rulecolor=\color{black},         % if not set, the frame-color may be changed on line-breaks within not-black text (e.g. comments (green here))
}

\usepackage[acronym, toc]{glossaries}
\makenoidxglossaries

\newacronym{cicd}{CI/CD}{Continuous Integration/Continuous Deployment}
\newacronym{http}{HTTP}{Hypertext Transfer Protocol}
\newacronym{api}{API}{Application Programming Interface}
\newacronym{crud}{CRUD}{Create, Read, Update, Delete}
\newacronym{erd}{ERD}{Entity-Relation-Diagram}
\newacronym{spa}{SPA}{Single-Page-Application}
\newacronym{css}{CSS}{Cascading Style Sheets}
\newacronym{tcp}{TCP}{Transmission Control Protocol}

\newglossaryentry{gitlabrunner}{name={GitLab-Runner},description={Anwendung, die auf Aktionen eines Git Repositories reagiert und vorkonfigurierte Jobs ausführt}}
\newglossaryentry{rest}{name={REST}, description = {Representational State Transfer, Ein Paradigma für die Architektur von \gls{crud} Webanwendungen }}
% ------------- Abbildungen mit Tikz erstellen
\usepackage{tikz}
\usetikzlibrary{arrows,automata,shapes,calc}


% ------------- Zeilenabstand
\usepackage{setspace}
\spacing{1.5}

% ------------- Kopfzeile
\usepackage[automark]{scrlayer-scrpage}
\pagestyle{scrheadings}
\clearscrheadings
\ihead{}
%\ihead[\rightmark]{\pagemark}
\chead{}
%\chead{}{}
%\ohead{\pagemark}
%\ohead[\pagemark]{\rightmark}
\ofoot{\pagemark}
%\ofoot[\pagemark]{\rightmark}


\renewcommand{\contentsname}{Inhaltsverzeichnis}
\renewcommand{\listfigurename}{Abbildungsverzeichnis}
\renewcommand{\listtablename}{Tabellenverzeichnis}

%-------------- Ende allgemeine Angaben -----------------------------
\begin{document}

%-------------- Beginn Titelseite -----------------------------------
\hypersetup{pageanchor=false}
\thispagestyle{empty}

\vspace{1cm}
\begin{center}
	\vspace{3cm}
	{\large Abschlussprüfung Sommer 2024} \\
	\vspace{2cm}
	{\large Fachinformatiker Anwendungsentwicklung} \\
	\vspace{3cm}
	{\huge \textbf{Titel}}\\
	\vspace{1.5cm}
	{\large Dokumentation zur betrieblichen Projektarbeit}\\
	\vspace{1.5pt}
	{\large IHK Hamburg }
\end{center}

\vspace{1cm}
\begin{flushleft}
	Abgabe: \\
	Hamburg, --.--.2024 
	
	\vspace{1cm}
	
	\textbf{Vorgelegt von:} \\
	\begin{tabular}{l l}
		Namen:       		& Nachname, Vorname \\
		Ausbildungsbetrieb: & ABC AG			\\
		Adresse:		    & Straße 12			\\
							& 22303 Hamburg		\\
	\end{tabular}
\end{flushleft}
%-------------- Ende Titelseite -------------------------------------

%-------------- Inhaltsverzeichnis
\newpage
\hypersetup{pageanchor=true}
\pagenumbering{arabic}
\renewcommand{\thepage}{\Roman{page}}
\tableofcontents{}

%-------------- Abbildungsverzeichnis
\newpage
\cleardoublepage\phantomsection
\addcontentsline{toc}{section}{Abbildungsverzeichnis}
\listoffigures
\renewcommand{\figurename}{Abb.}

%-------------- Abkürzungsverzeichnis
\newpage
\cleardoublepage\phantomsection
\printnoidxglossary[type=\acronymtype, title= Abkürzungsverzeichnis]

%-------------- Glossar
\newpage
\cleardoublepage\phantomsection
\printnoidxglossary[title= Glossar]

%-------------- Tabellenverzeichnis
\newpage
\cleardoublepage\phantomsection
\addcontentsline{toc}{section}{Tabellenverzeichnis}
\listoftables

\newpage
\pagenumbering{arabic}
\renewcommand{\thepage}{\arabic{page}}

%-------------- Einleitung
\section{Einleitung}

\subsection{Projektbeschreibung}

\subsection{Projektziel}

\subsection{Projektbegründung}

\subsection{Projektschnittstellen}

\subsection{Projektabgrenzung}

%------------- Projektplanung
\newpage
\section{Projektplanung}
%------------- Projektplanung

\subsection{Projektphasen}
Ein interessantes Buch zu dem Thema ist \cite{IOR2021}

\begin{table}[h]
	\begin{center}
		\rowcolors{2}{blue!80!green!50!yellow!10}{white}
		\begin{tabular}{l r}
			\rowcolor{blue!20}
			Projektphase 				  & Geplante Zeit \\
			\midrule
			AnalysePhase 			      &  5 h          \\
			Entwurfsphase   		      &  8 h          \\
			Implementierungsphase         & 50 h          \\
			Abnahme- und Deploymentsphase &  5 h		  \\
			Dokumentationsphase			  & 12 h		  \\
			\midrule
			Gesamt						  & 80 h		  \\
		\end{tabular}
	\end{center}
	\caption{Grobe Zeitplanung}
	\label{table:1}
\end{table}

\subsection{Ressourcenplanung}
%-------------- Analysephase
\newpage
\section{Analysephase}

\subsection{Ist-Analyse}

\subsection{Wirtschaftlichkeitsanalyse}

\subsubsection{Projektkosten}

\begin{table}[h]
	\begin{center}
		\rowcolors{1}{blue!80!green!50!yellow!10}{white}
		\begin{tabular}{l r r}
			Entwicklungszeit	&	   & 80 h   \\
			Gruppenmitglieder	&  4   &        \\
			Stundensatz      	& 15 € &        \\
			Ressourcenpauschale &  5 € &		\\
			Kosten Pro Stunde   &      & 65 €   \\
			\midrule
			\rowcolor{blue!20}
			Gesamt       		&  	   & 2500 € \\
		\end{tabular}
	\end{center}
	\caption{Projektkosten}
	\label{table:2}
\end{table}

\subsubsection{Nutzwertanalyse}

\begin{table}[h]
	\begin{center}
		\rowcolors{2}{blue!80!green!50!yellow!10}{white}
		\begin{tabular}{ l c | r r | r r}
											&				& \cellcolor{red!10!} Ist -	& \cellcolor{red!10!}	Zustand				& \cellcolor{green!10!} Projekt			& \cellcolor{green!10!}		\\
			\rowcolor{blue!10}
			Kriterien 						& \makecell{Gewicht\\in \%} & \makecell{Punkte\\(0-10)} & \makecell{Gewicht x\\Punkte} 	& \makecell{Punkte\\(0-10)} & \makecell{Gewicht x\\Punkte} \\
			\toprule
			Kosten       					& 20       		& 5   			& 100				& 3				&  60  \\
			Nachhaltigkeit 					& 30       		& 3   			&  90				& 7				& 210  \\
			Energieeffizienz 				& 40       		& 3   			& 120				& 7				& 280	\\
			\makecell{Automatisierungs-\\möglichkeiten}   & 10      		& 1     		&  10				& 8 			&  80	\\
			\hline
			\rowcolor{blue!20}
			Gesamt       					& 100      		& 12   			& 320				& 25 			& 630 	\\
		\end{tabular}
	\end{center}
	
	\caption{Nutzwertanalyse}
	\label{table:3}
\end{table}


%-------------- Entwurfsphase
\newpage
\section{Entwurfsphase}
%-------------- Entwurfsphase
\subsection{Architekturdesign}

\subsection{Entwurf der Benutzeroberfläche}\label{entwurfui}

\subsection{Datenmodell}

\subsection{Maßnahmen zur Qualitätssicherung}\label{quali}

\subsection{Deployment}

%-------------- Implementierungsphase
\newpage
\section{Implementierungsphase}
%-------------- Implementierungsphase

\subsection{Implementierung der Datenstrukturen}

\subsection{Implementierung Benutzeröberfläche}

\subsection{Implementierung der Middleware}
%-------------- Abnahme und Deployment
\newpage
\section{Abnahme und Deploymentsphase}

\subsection{Code-Review/ Team internes Review}

\subsection{Deplyoment}

%-------------- Fazit
\newpage
\section{Fazit}
%-------------- Fazit

\subsection{Soll-/Ist-Vergleich}

\begin{table}[h]
	\begin{center}
		\rowcolors{2}{blue!80!green!50!yellow!10}{white}
		\begin{tabular}{l r r r}	
			\rowcolor{blue!20}
			Projektphase 				 &Geplant& Tatsächlich & Differenz\\
			\midrule
			AnalysePhase 			      &  5 h &  7 h & + 2 h   \\
			Entwurfsphase   		      &  8 h &  6 h & - 2 h   \\
			Implementierungsphase         & 50 h & 52 h & + 2 h   \\
			Abnahme- und Deploymentsphase &  5 h &  5 h & - 1 h	  \\
			Dokumentationsphase			  & 12 h & 12 h & 	0 h   \\
			\midrule
			Gesamt						  & 80 h & 80 h &   0 h   \\
		\end{tabular}
	\end{center}
	\caption{Soll-/Ist-Vergleich}
	\label{table:4}
\end{table}

\subsection{Lessons Learned}

\subsection{Ausblick}

%-------------- Literaturverzeichnis
\newpage
\addcontentsline{toc}{section}{Literatur}
\bibliography{Quellen}

%-------------- Anhang
\newpage
\appendix
\section{Anhang}
\subsection{Detaillierte Zeitplanung}\label{A1}
\begin{table}[h]
	\begin{center}
		\rowcolors{2}{blue!80!green!50!yellow!10}{white}
		\begin{tabular}{l r r}	
			\rowcolor{blue!20}
			AnalysePhase 			       					&  		&  8 h  \\
			1. Durchführung einer detaillierten Ist-Analyse &  3 h 	&       \\
			2. Erstellung einer Wirtschaftlichkeitsanalyse 	&  2 h 	& 		\\
			3. Grundaufbau strukturieren 					&  3 h	& 		\\
			\rowcolor{blue!20}
			Entwurfsphase   		    					&  		&  4 h  \\
			1. Planung der Softwarearchitektur  			&  3 h  &	 	\\ 
			1. Planung der Bernutzeroberfläche  			&  1 h  &	 	\\ 
			\rowcolor{blue!20}
			Implementierungsphase    						&       & 50 h  \\	
			1. Frontend development 						& 15 h	& 		\\
			2. Backend development  						& 20 h  & 		\\
			3. Implementierung der Containerization  		& 15 h 	& 		\\
			\rowcolor{blue!20}
			Abnahme- und Deploymentsphase  					&      	&  6 h	\\
			1. Qualitätssicherung 							&  2 h 	& 		\\
			2. Interne Abnahme im Team  					&  2 h 	& 		\\
			3. Live-Stellung  								&  2 h 	& 		\\ 
			\rowcolor{blue!20}
			Dokumentationsphase			  					&     	& 12 h	\\
			1. Erstellung der Projektdokumentation  		& 12 h 	& 		\\
			\midrule
			\rowcolor{blue!20}
			Gesamt						   					&     	& 80 h	\\
		\end{tabular}
	\end{center}
\end{table}

\subsection{Ressourcen}\label{A2}
Hardware
\begin{itemize}
  \item Arbeitsplatz mit Laptop und externem Bildschirm für jeden Auszubildenden
  \item Debian 12 Linux Server
\end{itemize}

Software
\begin{itemize}
  \item Visual Studio Code - Entwicklungsumgebung
  \item XAMPP - Lokale MySQL Datenbank
  \item HeidiSQL - Grafischer SQL-Client
  \item Git - Verteilte Versionsverwaltung
  \item GitLab - Server für Git
  \item Vite - Buildtool
  \item Node.js - JavaScript Laufzeitumgebung
  \item MySQL - Datenbankmanagementsystem
  \item Vue.js - JavaScript Frontend Framework
  \item Docker - Software zur Containerization
  \item GitLab Runner - Pipeline
  \item Notion - Projektmanagement
  \item Discord - Kommunikationsplattform
  \item Google Docs - Kollaborative Notizplattform
\end{itemize}

Personal
\begin{itemize}
  \item Auszubildende - Umsetzung des Projektes
  \item Auszubildende - Review des Codes
  \item Mitarbeiter von Pemamek - Festlegung der Anforderungen, Abstimmung der Oberfläche und Abnahme des Projektes
\end{itemize}


\end{document}